\begin{lemma}[лемма к теореме Гильберта--Шмидта]\label{hilbert-smidt:lemma}
        Пусть $H(\C)$~--- гильбертово пространство, $A$~--- компактный самосопряжённый оператор, $A \neq 0$.
        Тогда у $A$ существует собственное значение, отличное от нуля.
\end{lemma}

\begin{proof}
    По теореме \nameref{11.5} $r(A)=max(|m_-, m_+|)= \|A\| \neq 0$ (так как оператор ненулевой).
    Так как $A \neq 0$, то хотя бы одно из этих чисел ($m_+, m_-$) не равно нулю. 
    То, которое не ноль, по лемме \nameref{fredholm:lemma1} это $m_* \neq 0$~--- собственное значение.
\end{proof}

\begin{theorem}[12.5 (Гильберт--Шмидт)]\label{12.5}
    $H(\C)$~--- сепарабельное гильбертово пространство, $A \in \mathcal{L}(H)$~--- компактный самосопряжённый оператор. Тогда в $H$ существует ортонормированный базис, состоящий из собственных векторов оператора $A$.
\end{theorem}

\begin{proof}
    Вспоминаем, что $A$~--- компактный оператор. По теореме \nameref{12.3} внешность любой окрестности нуля содержит конечное число собственных значений. Это позволит нам упорядочить ненулевые собственные значения по невозрастанию, взяв сначала положительные из равных мо модулю (\(|\lambda_1| \geqslant |\lambda_2| \geq \ldots \)).

    Поскольку для любого $\lambda_k$ $\dim \ker A_{\lambda_k} = m_k < \infty$ по \nameref{12.2}, в нём можно выбрать ОНБ и назвать его векторы $e^k = \left\{ e_i^k \right\}_{i=1}^{m_k}$. Заметим, что эти векторы будут собственными для $A$. Введём $E_1 = \bigcup\limits_{i=1}^\infty e^l$

    Далее рассмотрим $\ker A \subset H$, оно будет сепарабельным, значит в нём можно выбрать ОНБ $\{ E_0 \}$. Обозначим $e = E_0 \cup E_1$.  Все элементы $\{e\}$ ортогональны (по теореме \nameref{11.1}, так как собственные векторы соответствующие разным собственным значениям ортогональны).

    $\{e\}$~--- ортонормированная система. Как доказать, что она базис?

    Вспоминаем: ортонормированная система в $H$ является базисом $\lra M:=\overline{[\{e\}]}=H$.
    Запишем теорему Рисса о проекции: $M \oplus M^\perp = H$. Нужно доказать, что $M^\perp = \{0\}$.

    $M$ инвариантно относительно $A$: если $x \in M$, то $x = \lim\limits_{n \rightarrow \infty} x_n =  \lim \limits_{n \rightarrow \infty} \sum \limits_k \alpha_{n_k} e_k$ для $\left\{ x_n \right\} \subset [\left\{e \right\}]$. Подействовав оператором $A$, получаем $Ax = \lim \limits_{n \rightarrow \infty} \sum \limits_k \alpha_{n_k} A e_k = \lim \limits_{n \rightarrow \infty} \sum\limits_{k}\alpha_{n_k}\lambda_k e_k$.

    По лемме \nameref{fredholm:lemma3} $M^\perp$ инвариантен относительно $A$.
    Рассмотрим оператор $\tilde{A} = \left. A \right|_{M^\perp}$
    ~--- компактный самосопряжённый оператор. Рассмотрим два варианта:
    \begin{itemize}
        \item[\(\tilde{A} = 0\)]
        Это означает, что $\tilde{A} x = 0 \: \forall x \in M^\perp$. Следовательно, 
        $M^\perp \subseteq \ker \tilde{A}$. Значит, если $M^\perp \neq \left\{0 \right\}$, то должны быть собственные векторы $\tilde{A}$ из него, а это противоречит тому, что все собственные векторы в $M$. Значит, $M^\perp = \left\{0 \right\}$
        \item[\(\tilde{A} \neq 0\)]
        Тогда, по \hyperref[hilbert-smidt:lemma]{лемме к Т. Гильберта--Шмидта} у $\tilde{A}$ есть собственные значение $\lambda \neq 0$ (оно же является собственным значением оператора $A$). У $\lambda$ есть соответсвующий собственный вектор, а это входит в противоречие с тем, что все собственные векторы в $M$. Значит, $M^\perp = \left\{0 \right\}$.
    \end{itemize}
\end{proof}